%!TEX root = main.tex
\section{Related Work}
\label{sec:related_work}

\textbf{Vietnamese hate speech.}  The UIT-ViHSD corpus \cite{luu2021vihsd}
established the three-class \hate{}/\off{}/\clean{} task.  Prior work has
explored monolingual PhoBERT \cite{nguyen2020phobert}, multilingual mBERT/XLM-R
\cite{conneau2020xlmr}, and data augmentation \cite{nguyen2026vicm}.
Ensemble and curriculum learning \cite{bengio2009curriculum} raised macro-F1 to
0.66--0.68, but the decision process remains a black box.

\textbf{English hate speech.}  The Davidson et al. \cite{davidson2017automated}
dataset (Hate Speech and Offensive Language) is the de facto benchmark for
English \hate{}/\off{}/\clean{}.  Fine-tuned BERT/RoBERTa reach 0.77--0.81
macro-F1; recent work uses multi-task learning and adversarial training
\cite{macavaney2019hate}.

\textbf{Hybrid rule + neural systems.}  Early hybrids concatenated lexical
features to neural inputs \cite{waseem2016hateful}.  \cite{vidgen2021directions}
argue for interpretable components.  \cite{bojchevski2020fusing} fuse rule
scores with neural logits via learnable weights; \cite{ahn2022evidence}
propose evidence channels.  Most approaches evaluate on a single language and
do not provide span-level evidence or language-aware routing.

\textbf{Language-aware routing.}  Mixed-language moderation is addressed by
language identification \cite{lui2012langid} followed by per-language models
\cite{pires2019multilingual}.  Code-mixed text (Vi-En) is handled by
multilingual backbones or dedicated tokenizers \cite{nguyen2026vicm}.  Our
design differs by keeping both rule engines active on \texttt{mix} and
fallback segments, and by fusing per-language scores rather than training a
single multilingual head.

\textbf{Explainable moderation.}  Post-hoc explanations (LIME, SHAP,
attention) are common \cite{danilevsky2020survey}.  Our approach provides
\emph{intrinsic} evidence: matched lexicon spans, match types and arithmetic
scores, which survive deployment without additional computation.